\section{Hardness under the Small Set Expansion Hypothesis}
\label{sec:sseh}

We prove that the Small Set Expansion Hypothesis rules out every fixed
approximation factor for treedepth. The source is the global partition
promise of Raghavendra, Steurer, and Tulsiani~\cite{RST}. We call
the source graph the \emph{host} to distinguish it from its line
graph. For a fixed number $q$ of classes, the YES case has a partition
into classes of nearly equal volume, each with few crossing edges.
For fixed source parameters $c,\varepsilon>0$, the NO case gives every
set of normalized volume in $[1/4,3/4]$ at least
$(c/2)\sqrt{\varepsilon}\,m$ boundary edges. Passing to
the line graph turns the YES partition into an elimination forest of
height $O(\varepsilon m)$ and the NO cut bound into a balanced-separator
lower bound of $\Omega(\sqrt{\varepsilon}\,m)$, where $m$ is the number
of host edges and hence the order of the line graph.

% Wu, Austrin, Pitassi, and Liu's incidence construction also gives
% SSE hardness for treedepth at every fixed factor when combined with
% the elimination-forest bound in Appendix~\ref{app:incidence-proof}. Their output has one
% part for replicated host vertices and another for host edges.

The line-graph construction uses host edges as vertices, joining two
when the corresponding edges share an endpoint. This gives separator
bounds directly in terms of the output order, which we use again in
the bounded-degree reductions. We use an unweighted host promise.
Appendix~\ref{app:sse-normalization} derives it from the weighted
source by rounding edge weights.

\subsection{The source partition promise}

For a finite undirected weighted graph $G$, let $W>0$ be the total
undirected edge weight. Edge weights are positive rationals, and
parallel edges and loops are allowed. For a multiset $F$ of edges,
$w(F)$ denotes its total weight, counting multiplicity.
The weighted degree $\deg_G(v)$ is the sum of
the weights of the edges incident to $v$, with a loop counted twice.
Define $\vol_G(A)=\sum_{v\in A}\deg_G(v)$ and
$\mu_G(A)=\vol_G(A)/(2W)$. When $\vol_G(A)>0$, put
$\Phi_G(A)=w(\cut_G(A))/\vol_G(A)$. Thus a loop contributes twice
its weight to the volume and never contributes to a cut.

\begin{lemma}[Global SSE partition gap, normalized form]\label{lem:sse-unweighted-gap}
There are absolute constants $c>0$ and $\varepsilon_0>0$ such that,
for every fixed integer $q\geq2$ and every fixed rational
$0<\varepsilon<\varepsilon_0$, the following promise on unweighted
multigraphs $G$ with $m>0$ edges, allowing loops, is SSE-hard:
\begin{description}
\item[YES:] There is a partition
  $V(G)=V_1\mathbin{\dot\cup}\cdots\mathbin{\dot\cup}V_q$ such that
  $(1-\varepsilon)/q\leq\mu_G(V_i)\leq(1+\varepsilon)/q$ and
  $\Phi_G(V_i)\leq3\varepsilon$ for every $i$.
\item[NO:] Every $A\subseteq V(G)$ with
  $\mu_G(A)\in[1/4,3/4]$ satisfies
  $|\cut_G(A)|\geq(c/2)\sqrt{\varepsilon}\,m$.
\end{description}
\end{lemma}

Appendix~\ref{app:sse-normalization} derives this promise from the
published weighted gap by rounding edge weights to parallel edge
multiplicities. The host promise imposes no maximum-degree bound.
The later degree reduction acts on the line graph through incidence
clique factors, in Section~\ref{subsec:bounded-degree-consequences}.
The following counting lemma converts the partition properties into
bounds on crossing and internal edge weights.

\begin{lemma}[Crossing and internal edge weights]
\label{lem:sse-yes-edge-counts}
Let $G$ be a finite undirected weighted multigraph with positive
rational edge weights and total edge weight $W>0$.
Let $(V_1,\ldots,V_q)$ be a partition of its vertices into parts
of positive volume.
Let $D$ be the multiset of edges whose endpoints lie in different
classes. For each $i$, let $E_i$ be the multiset of edges with both
endpoints in $V_i$, including loops at vertices of $V_i$. Each
parallel edge is included in the corresponding multiset. Then
$w(D)\leq W\max_i\Phi_G(V_i)$ and
$w(E_i)\leq\vol_G(V_i)/2$ for every $i$.
\end{lemma}

\begin{proof}
Each edge in $D$ belongs to the cuts of its two endpoint classes.
Internal edges, including loops, belong to none. Hence
\[
  2w(D)=\sum_{i=1}^q w(\cut_G(V_i))
  =\sum_{i=1}^q\Phi_G(V_i)\vol_G(V_i)
  \leq\max_i\Phi_G(V_i)\sum_{i=1}^q\vol_G(V_i)
  =2W\max_i\Phi_G(V_i).
\]
Each internal edge, including a loop, contributes twice its weight
to its class volume, whereas a crossing edge contributes once. Thus
$2w(E_i)+w(\cut_G(V_i))=\vol_G(V_i)$, which gives the second bound.
\end{proof}

\subsection{The line graph and its separator lower bound}

For an unweighted host multigraph $G$, its \emph{line graph}
$\mathcal{L}(G)$ has one vertex for each edge of $G$. Parallel edges
give distinct vertices. Two distinct output vertices are adjacent
if and only if their host edges share an endpoint.
This is a simple unweighted graph, even when the host has
parallel edges or loops. If the host has $m$ edges, counting parallel
edges separately, then $\mathcal{L}(G)$ has $m$ vertices.

\begin{lemma}[A treedepth upper bound from a partition]\label{lem:sse-complete}
Let $G$ be an unweighted host multigraph with $m>0$ edges, counting
parallel edges separately, and let $(V_1,\ldots,V_q)$ be a partition
of its vertices. If $D$ is the multiset of edges joining different
classes, then $\td(\mathcal{L}(G))\leq |D|+\max_i |E(G[V_i])|$.
\end{lemma}

\begin{proof}
Make one chain from the line-graph vertices corresponding to the
edges in $D$. For each class $V_i$, make a second chain from the
vertices corresponding to the internal edges
of $G[V_i]$. If $D$ is nonempty, attach the class chains below the last
vertex of the $D$-chain. If $D$ is empty, use the class chains as
separate roots.

Vertices on the $D$-chain are comparable with all vertices of the
forest. Within each class chain all vertices are comparable as well.
The only remaining pairs lie on different class chains. Their host
edges have endpoints in disjoint classes, so those pairs are
nonadjacent in $\mathcal L(G)$. Thus the constructed forest is an
elimination forest of height $|D|+\max_i|E(G[V_i])|$, including
when $D$ is empty.
\end{proof}

For the converse direction, deleting vertices of $\mathcal L(G)$
means deleting the corresponding multiset of host edges. Different
components of the remaining line graph have disjoint sets of incident
host vertices. We use these sets to recover a host cut whose edges
must all have been deleted.

For $0<\rho<1$, a \emph{$\rho$-balanced separator} of an $N$-vertex
graph $H$ is a set $S\subseteq V(H)$ such that every component of
$H-S$ has at most $\rho N$ vertices. We write $b_\rho(H)$ for its
minimum size. A \emph{balanced separator} is a $1/2$-balanced
separator, and we abbreviate $b_{1/2}(H)$ to $b(H)$.

\begin{lemma}[Balanced separators from tree decompositions
{\cite[Lemma~5.3]{APW}}]
\label{lem:treewidth-balanced-separator}
Every finite nonempty graph $H$ satisfies
$b(H)\leq\operatorname{tw}(H)+1$.
\end{lemma}

\begin{lemma}[From host cuts to line-graph separators]
\label{lem:linegraph-separators}
Let $G$ be an unweighted host multigraph with $m>0$ edges, counting
parallel edges separately.
Suppose that, for some $\beta>0$, every set $A$ with
$\mu_G(A)\in[1/4,3/4]$ satisfies $|\cut_G(A)|\geq\beta m$.
Then $b_\rho(\mathcal{L}(G))\geq\min\{(3/4-\rho)m,\beta m\}$
for every $1/2\leq\rho<3/4$.
\end{lemma}

\begin{proof}
Fix a $\rho$-balanced separator $S$ of $\mathcal L(G)$, and write
$F_S$ for the corresponding multiset of host edges, so $|F_S|=|S|$.
If $|S|\ge(3/4-\rho)m$, the claim is immediate. Assume henceforth
that $|S|<(3/4-\rho)m$.

\emph{Grouping the surviving host edges.}
The line graph $\mathcal L(G)-S$ has
$m-|S|>(1/4+\rho)m$ vertices, with each component of order at
most $\rho m$. Choose a union $Z$ of components with
$m/4\le|Z|\le\rho m$. If a component has size in this interval,
take it. Otherwise, every component has fewer than $m/4$ vertices.
Add components until their total first reaches $m/4$.
The total is then below $m/2\le\rho m$.

Let $X$ be the host edges corresponding to vertices in $Z$, and
let $Y$ be the other surviving host edges. Then
$|X|=|Z|\ge m/4$ and
$|Y|=m-|S|-|X|>m/4$.
If a host vertex were incident to edges in both multisets, their
line-graph vertices would be adjacent, with one in $Z$ and one
outside $Z$. This contradicts the choice of $Z$ as a union of
components.

\emph{Forming the host cut.}
Let $A,B\subseteq V(G)$ be the sets of endpoints of edges in
$X,Y$, respectively. The preceding paragraph gives
$A\cap B=\varnothing$. Every edge of $X$
contributes two degree units to $\vol_G(A)$, including a loop.
The same holds for $Y$ and $B$. Thus
$\vol_G(A)\ge2|X|\ge m/2$ and
$\vol_G(B)\ge2|Y|\ge m/2$.
Since the total host volume is $2m$,
$\vol_G(A)\le2m-\vol_G(B)\le3m/2$.
Since $\mu_G(A)=\vol_G(A)/(2m)$, we obtain
$\mu_G(A)\in[1/4,3/4]$.
Every surviving edge lies wholly in $A$ or wholly in
$B\subseteq V(G)\setminus A$, so every edge crossing $A$ belongs
to $F_S$. The cut assumption therefore yields
$|S|=|F_S|\ge|\cut_G(A)|\ge\beta m$, proving the bound.
\end{proof}

\begin{lemma}\label{lem:balanced-by-td}
For every finite graph $H$, $b(H)\leq\td(H)$.
Moreover, every elimination forest $F$ of an $N$-vertex graph $H$
satisfies $\totd(F)\geq(N/2)b(H)$.
\end{lemma}

\begin{proof}
Fix an elimination forest $F$ of $H$. If every root subtree has
at most $N/2$ vertices, then $b(H)=0$ and both claims follow.
Otherwise, delete the root of the unique subtree of size greater
than $N/2$, and continue into a child subtree of size greater than
$N/2$ while one exists. Such a child is unique, since two disjoint
subtrees of that size cannot fit among $N$ vertices.

The deleted vertices form an ancestor path $P$.
Every remaining
subtree has at most $N/2$ vertices, and no graph edge joins distinct
remaining subtrees. Thus $P$ is balanced and
$b(H)\leq|P|\leq h(F)$. Minimizing height proves the first claim.
The original subtree rooted at the last vertex of $P$ has more than
$N/2$ vertices, each at depth at least $|P|\geq b(H)$. Their depths
alone give $\totd(F)\geq(N/2)b(H)$.
\end{proof}

\begin{corollary}[NO treedepth bound]\label{lem:sse-sound}
Under the cut assumption of Lemma~\ref{lem:linegraph-separators},
$\td(\mathcal L(G))\geq\min\{m/4,\beta m\}$.
\end{corollary}

\begin{proof}
Take $\rho=1/2$ in Lemma~\ref{lem:linegraph-separators} and use
$b(\mathcal L(G))\leq\td(\mathcal L(G))$ from
Lemma~\ref{lem:balanced-by-td}.
\end{proof}

\subsection{The approximation gap}
\label{subsec:sse-dense-gap}

\begin{proof}[Dense treedepth gap]
The YES and NO bounds below have coefficients $4\varepsilon$ and
$(c/2)\sqrt\varepsilon$. Their ratio $c/(8\sqrt\varepsilon)$ grows
without bound as $\varepsilon$ decreases, so a sufficiently small
fixed source error gives any desired constant gap.
Fix $C\geq1$ and choose a fixed rational $\varepsilon>0$ small enough
that Lemma~\ref{lem:sse-unweighted-gap} applies,
$c\sqrt\varepsilon\leq1/4$, and $\sqrt\varepsilon<c/(8C)$.
Since $\varepsilon$ is fixed, $q=\lceil2/\varepsilon\rceil$ is a
constant. Apply that lemma with this $q$ to obtain an
unweighted host $G^*$ with $m$ edges, and output $\mathcal L(G^*)$.

In the YES case, let $D$ be the multiset of host edges joining different
partition classes. Lemma~\ref{lem:sse-yes-edge-counts} gives
$|D|\leq3\varepsilon m$ and
$|E(G^*[V_i])|\leq\vol_{G^*}(V_i)/2\leq(1+\varepsilon)m/q$.
Lemma~\ref{lem:sse-complete} therefore gives
\[
\begin{aligned}
 \td(\mathcal L(G^*))
 &\leq3\varepsilon m+(1+\varepsilon)m/q
 \leq3\varepsilon m+(1+\varepsilon)\varepsilon m/2\\
 &=(7/2+\varepsilon/2)\varepsilon m<4\varepsilon m,
\end{aligned}
\]
where $q=\lceil2/\varepsilon\rceil$ gives $1/q\leq\varepsilon/2$,
and $\varepsilon<1$ gives $7/2+\varepsilon/2<4$.
In the NO case, the choice $c\sqrt{\varepsilon}\leq1/4$ makes
$(c/2)\sqrt{\varepsilon}\leq1/8<1/4$. Thus the smaller term in
the minimum from Corollary~\ref{lem:sse-sound} is the cut bound, and
that corollary gives
\[
 \td(\mathcal L(G^*))\geq(c/2)\sqrt\varepsilon\,m
 >4C\varepsilon m,
\]
since $\sqrt\varepsilon<c/(8C)$ implies
$(c/2)\sqrt\varepsilon>4C\varepsilon$. Thus the NO bound exceeds the YES
bound by a factor greater than $C$.
\end{proof}
